\section{Use structure to show priority}
\label{sec:structure}

A document can look organized without helping anyone find the main point.
Agent-like answers often add headings and bullets before mapping the
dependencies among ideas. The result gives every item the same visual weight.

\subsection{Use lists for coherent sets}

A list helps when readers need to scan a set or follow a sequence. Procedures
often qualify; fixed option sets do too. Connected claims usually belong in
prose because their relation matters more than their count.

A useful editing heuristic is to avoid introducing lists longer than five
items unless the complete set matters. This is not a mathematical limit. It
forces the writer to distinguish one coherent set from several unnamed groups.

Do not force material into threes. Three claims should appear together because
the subject has three relevant parts, not because a triad sounds finished.
The same rule applies to adjectives and examples. Avoid false ranges for the
same reason.
``From hiring to reporting'' implies a scale that does not exist; ``staffing
and reporting'' simply names the topics.

\subsection{Use headings for navigation}

A heading should mark a real change in question or task. Remove headings over
one-line sections and sentences that merely repeat the heading. Bold inline
labels often make ordinary prose look like a dashboard.

Announcements such as ``Let's dive in,'' ``Here's what you need to know,'' and
``Now we will explore'' tell the reader that content is coming. Start with the
content. Similarly, documentation should describe the current system rather
than narrate the latest diff unless the document is a changelog, release note,
or migration guide.

\subsection{Vary rhythm without forcing drama}

Natural prose changes pace. It may use a short sentence for emphasis, but it
does not make every line a punchline. Combine stacked fragments when they
manufacture urgency:

\agentdraft

Then the new model arrived. No warning. No compromise. The old rules were
gone.

\humanrevision

The new model changed the comparison because it did not share the assumptions
of the earlier systems.

The revision gives the reader a reason for the change. It does not rely on
cinematic timing to make the change feel important.

Avoid slogan formulas for the same reason. ``Symmetry is the language of
trust'' sounds memorable but leaves the relation unspecified. If symmetric
layouts make controls easier to predict, state that claim and its conditions.

\subsection{Formatting is weak evidence}

Mechanical boldface and title case in every heading often accompany agent-like
prose. Decorative emoji and inline-header lists do too. Remove these devices
when they obscure hierarchy or clash with the publication. Do not treat them
as proof of authorship.

Punctuation also follows the venue and the writer. The Humanizer pass removes
dashes because agents often use them as a default rhythm device. The broader
general-writing standard allows an occasional dash in a long piece when it
fits the writer's sample better than the alternatives. A comma may work, or
the sentence may need a colon or parenthesis. Curly quotation marks are
likewise a typographic convention, not a reliable signal. Match the
document's style rather than changing punctuation for forensic reasons.

\subsection{Repeat stable terms}

Synonym cycling makes a stable object appear to change. If a paragraph begins
with ``the model,'' do not rotate through ``the system,'' ``the framework,''
and ``the approach'' merely to avoid repetition. Repeat the clear term until a
real distinction requires another one. Controlled repetition reduces the
reader's memory burden.
